% !TeX root = ../nc-cubic-unmark.tex
\section{数値的線束類の実現と弱標識}\label{sec:line-classification}
固定した Abel 圏の弱標識忘却を理解するためには、まずその圏の中の
線束を分類し、次に exact 自己同値による軌道を求める必要がある。
本節では rank-one torsion-free exceptional 対象を\emph{線束対象}と呼ぶ。
その同型類集合を $\operatorname{Line}(\mathcal A)$、
豊富性と式\eqref{eq:weak-hilbert}も満たす部分を
$\operatorname{Line}_{\rm adm}(\mathcal A)$ と書く。
対象同士のテンソル積も、この集合上の群法則も仮定しない。

以下の主張は固定した反標準 heart の公理から述べる。
数値的線束類の存在は、一意性とは別に証明する必要がある。
内在的線束判定の新規性については網羅的な優先権確認を行っていない。

\subsection{例外性は自己 Euler 形式で判定できる}
以下、定義\ref{def:weak-cubic}を満たす $(\mathcal A,\OO_0)$ を固定する。
$H$ は exact なので $\mathcal A_1$ と rank-one torsion-free 性を保つ。

\begin{lemma}[階数一の射]\label{lem:rank-one-monomorphism}
rank-one torsion-free 対象 $E,F$ の間の非零射は単射である。
また $\End(E)=\mathbb C$ である。
\end{lemma}
\begin{proof}
非零射 $f:E\to F$ に対して $qf=0$ なら、$\operatorname{im}f$ は
$\mathcal A_1$ に属して $F$ に埋め込まれ、torsion-free 性に矛盾する。
従って $qf$ は非零単純対象間の同型である。
$\ker f\in\mathcal A_1$ が $E$ に埋め込まれるので $\ker f=0$ となる。
同じ理由で $\End(E)\to\End(qE)$ は単射である。
右辺は斜体、左辺は Hom 有限性により有限次元の複素代数であるから、
左辺は零因子を持たない有限次元代数である。
任意の元の最小多項式は $\mathbb C$ 上で一次式に分解し、
零因子がないことからその元はスカラーとなる。
\end{proof}

\begin{theorem}[内在的な線束判定]\label{thm:intrinsic-line-criterion}
任意の rank-one torsion-free 対象 $E$ と整数 $b>0$ に対して
\[
 \Hom(E,H^{-b}E)=0,\qquad \Ext^2(E,E)=0,
 \qquad \dim\Ext^1(E,E)=1-\chi(E,E).
\]
特に
\begin{equation}\label{eq:intrinsic-line-criterion}
 E\text{ が線束対象}\quad\Longleftrightarrow\quad\chi(E,E)=1.
\end{equation}
この主張には $E$ 自身の Artin--Zhang 豊富性を仮定しない。
\end{theorem}
\begin{proof}
$\OO_0$ の豊富性から有限直和による全射
$\bigoplus_i H^{-a_i}\OO_0\twoheadrightarrow E$ がある。
非零な成分を選ぶと、前補題により $H^{-a}\OO_0\hookrightarrow E$ を得る。
$h_E(n)=\dim\Hom(\OO_0,H^nE)$ と置けば、$n\ge a$ に対し
\[
 h_E(n)\ge 1+\tfrac32(n-a)(n-a+1).
\]
もし $E\to H^{-b}E$ が非零なら、それも単射であり、$H$ の exact 性から
$h_E(n)\le h_E(n-b)$ がすべての $n$ で成り立つ。
各剰余類 $n\bmod b$ で $h_E(n)$ は上に有界となり、上の二次増大と矛盾する。
Serre 双対性から
$\Ext^2(E,E)=\Hom(E,H^{-1}E)^*=0$ である。
補題\ref{lem:dimension-centre}によって高次 Ext は消え、
前補題の $\dim\End(E)=1$ と Euler 形式の定義から結論を得る。
\end{proof}

\begin{remark}[rank-one だけでは足りない]
可換な滑らかな三次曲面では $E=L\otimes I_Z$ に対し
$\chi(E,E)=1-2\operatorname{length}Z$、
$\dim\Ext^1(E,E)=2\operatorname{length}Z$ となる。
特に点のイデアル層 $I_p$ は rank-one torsion-free だが例外的でない。
一方、定義\ref{def:weak-cubic}の他の条件を保持するなら、
例外性の条件は $\chi(\OO,\OO)=1$ に置き換えてよい。
上の証明が使う $\OO_0$ の条件は rank-one、豊富性と切断数の増大であり、
この置換の確認に例外性を先取りする必要はない。
\end{remark}

\subsection{数値的欠損と同一数値類での一意性}
次の三次曲面型の数値的構造は\emph{追加仮定}である。
$K_0^{\rm num}(\mathcal A)$ に三次曲面型の表示
$\operatorname{ch}(E)=(r,D,s)$、$r\in\mathbb Z$、$D\in\Lambda=I_{1,6}$、
$K^2=3$ を指定し、$s-D^2/2\in\mathbb Z$、
\begin{align}
 \chi((r,D,s),(r',D',s'))
 &=rr'+rs'+r's-D\cdot D'
   -\tfrac12 K\cdot(rD'-r'D),\label{eq:line-rr}\\
 \operatorname{ch}(HE)
 &=(r,D-rK,s-D\cdot K+\tfrac12rK^2)\label{eq:line-H}
\end{align}
をすべての対象について仮定し、$r(\OO_0)=1$ と正規化する。
数値的階数と内在的 rank 条件の対応、および torsion 対象の次数正値性は、
以下でこの仮定から導く。
数値的線束類とは $v_D=(1,D,D^2/2)$ のことである。

\begin{theorem}[反標準 Serre 消滅と内在的階数]\label{thm:ac-serre-positivity}
定義\ref{def:weak-cubic}と式\eqref{eq:line-rr}--\eqref{eq:line-H}、
$r(\OO_0)=1$ を仮定する。$d(M)=(-K)\cdot c_1(M)$、$d_0=d(\OO_0)$ と置く。
任意の $M\in\mathcal A$ に対して
\begin{equation}\label{eq:ac-serre-vanishing}
 \Ext^i(\OO_0,H^nM)=0\qquad(i>0,\ n\gg0)
\end{equation}
であり、十分大きな $n$ で
\begin{equation}\label{eq:ac-hilbert-polynomial}
 \dim\Hom(\OO_0,H^nM)
 =\tfrac32r(M)n(n+1)+(d(M)-r(M)d_0)n+\chi(\OO_0,M).
\end{equation}
さらに次が成り立つ。
\begin{enumerate}
\item 全対象で $r(M)\ge0$。$r(M)=0$ なら $d(M)\ge0$。
\item $r(M)=d(M)=0$ なら $M\in\mathcal A_0$。$M\ne0$ なら
      $\chi(\OO_0,M)>0$ であり、$\operatorname{length}M\le\chi(\OO_0,M)$。
\item $\mathcal A_1=\{M:r(M)=0\}$。rank-one torsion-free 対象 $E$ には $r(E)=1$。
\end{enumerate}
特に有効次数正値性
\begin{equation}\label{eq:effective-degree-positivity}
 d(T)\ge0\qquad(T\in\mathcal A_1) \tag{EP}
\end{equation}
は独立の仮定を要しない。
\end{theorem}
\begin{proof}
$D_0=c_1(\OO_0)$ とする。$\OO_0$ の例外性と Riemann--Roch より
$\operatorname{ch}(\OO_0)=v_{D_0}$ であり、
$\chi(\OO_0,H^m\OO_0)=1+3m(m+1)/2$ である。
$m\ge0$ では Serre 双対性と所定の負の Hom 消滅から
$\Ext^2(\OO_0,H^m\OO_0)=0$ である。
式\eqref{eq:weak-hilbert}との差を取れば
$\Ext^1(\OO_0,H^m\OO_0)=0$ も得る。

豊富性により $P=\bigoplus_jH^{-a_j}\OO_0\twoheadrightarrow M$ を選ぶ。
この全射は
\[
 \Hom(M,H^{-n-1}\OO_0)
 \hookrightarrow\bigoplus_j\Hom(\OO_0,H^{a_j-n-1}\OO_0)
\]
を与える。右辺は $n\gg0$ で零なので、Serre 双対性から
$\Ext^2(\OO_0,H^nM)=0$ となる。
$N=\ker(P\to M)$ にもこの第二 Ext 消滅を適用すると、
長完全列の
\[
 \Ext^1(\OO_0,H^nP)\longrightarrow\Ext^1(\OO_0,H^nM)
 \longrightarrow\Ext^2(\OO_0,H^nN)
\]
の両端が $n\gg0$ で零となる。
高次 Ext は補題\ref{lem:dimension-centre}で消える。
これで式\eqref{eq:ac-serre-vanishing}を得て、Riemann--Roch を展開すれば
式\eqref{eq:ac-hilbert-polynomial}となる。

非零対象 $M$ ではこの多項式は恒等的に零でない。
実際、AZ 豊富性は任意に負の tail 内の有限直和による生成を与える
\cite[Lemma 2.1(ii)]{PolishchukProj}ので、
$\Hom(\OO_0,H^nM)$ は任意に大きい $n$ のうち少なくとも一つで非零である。
非負の値を取る非零多項式の最高次係数を順に見ると、(1)と、
$r(M)=d(M)=0$, $M\ne0$ のときの $\chi(\OO_0,M)\in\mathbb Z_{>0}$ を得る。

$r(M)=d(M)=0$ のとき、加法性と (1) により全部分商の $r,d$ も零である。
従って任意の有限な厳密部分対象列の各非零商に対して
$\chi(\OO_0,-)$ は正整数であり、列の長さは $\chi(\OO_0,M)$ 以下である。
これが (2) を示す。

次に $r(M)=0$ とし、$Q_0:\mathcal A\to\mathcal A/\mathcal A_0$ と置く。
$Q_0M$ の有限な厳密部分対象列は、Serre 商の部分対象の持上げにより
$M$ 内の有限な部分対象列で表せる。必要なら代表の累積和を取り、
始点を $0$、終点を $M$ とする。
各厳密段階の商 $N$ は $r(N)=0$ かつ $Q_0N\ne0$ だから、(2) より $d(N)>0$ である。
$d(N)$ は整数なので、この列の長さは $d(M)$ 以下である。
従って $Q_0M$ は有限長、すなわち $M\in\mathcal A_1$ となる。
この議論は $d$ が $\mathcal A_0$ 上で零になることを用いていない。

逆に $T\in\mathcal A_1$ に対し、全射
$\bigoplus_jH^{-a_j}\OO_0\twoheadrightarrow T$ の各成分の像を $I_j$ とする。
$\mathcal A_1$ は Serre 部分圏なので $I_j\in\mathcal A_1$ であり、
(1)と $r(\OO_0)=1$ から $r(I_j)\in\{0,1\}$ である。
$r(I_j)=1$ なら、その成分の核の階数は零、従って既に示した向きにより
核も $\mathcal A_1$ に入る。すると $H^{-a_j}\OO_0\in\mathcal A_1$ となり、
$q\OO_0\ne0$ と矛盾する。
従って全 $I_j$ の階数は零であり、その有限和 $T$ も階数零である。
これで $\mathcal A_1=\{r=0\}$ を得る。

rank-one torsion-free 対象 $E$ には豊富性と
補題\ref{lem:rank-one-monomorphism}から $H^{-a}\OO_0\hookrightarrow E$ が存在する。
商圏で非零単純対象間の同型になるので、その余核は $\mathcal A_1$ に入る。
従って $r(E)=1$。最後に (EP) は (1)と (3) の帰結である。
\end{proof}

\begin{remark}[入力と結論の範囲]\label{rem:serre-inputs}
上の消滅は一般の AZ 豊富性だけから導いたものではない。
指定した Serre 関手、三次曲面型 Riemann--Roch、構造対象の Hilbert 関数を使っている。
生成には異なる指数の有限直和だけを使い、各 $n\gg0$ について
単一の $H^{-n}\OO_0$ の直和で生成できるとは主張しない。
また (2) の逆、すなわち $M\in\mathcal A_0\Rightarrow d(M)=0$ もここでは未証明である。
\end{remark}

\begin{corollary}[零次元欠損の数値的類似]\label{cor:line-numerical-defect}
rank-one torsion-free 対象 $E$ の表示を $(1,D,s)$ とし、
$n(E)=D^2/2-s$ と置く。このとき
\[
 n(E)\in\mathbb Z_{\ge0},\qquad
 \chi(E,E)=1-2n(E),\qquad
 \dim\Ext^1(E,E)=2n(E).
\]
従って数値的線束類を持つ rank-one torsion-free 対象は自動的に例外的である。
\end{corollary}
\begin{proof}
式\eqref{eq:line-rr}に $E$ を二度代入し、
定理\ref{thm:intrinsic-line-criterion}を適用する。
\end{proof}

一意性には自己 Ext だけではなく二対象間の Ext が必要である。
以下の Serre 双対性と次数による消滅論法は、異なる設定である
terminal del Pezzo order 上の先行例\cite[Lemma 3.9]{HR}とも共通する。
そのため定理\ref{thm:ac-serre-positivity}で導いた (EP) を用いる。
有限代数の Euler 行列だけの一致は、この定理の反標準 heart の仮定を代用しない。

\begin{theorem}[数値的線束類における一意性]\label{thm:numerical-line-uniqueness}
定理\ref{thm:ac-serre-positivity}の数値仮定を置く。
rank-one torsion-free 対象 $E,F$ が $d(E)<d(F)+3$ を満たせば
$\Ext^2(E,F)=0$ である。特に $[E]=[F]=v_D$ なら $E\cong F$ である。
従って $\operatorname{Line}(\mathcal A)$ から $\Lambda$ への第一 Chern 類写像は単射である。
\end{theorem}
\begin{proof}
Serre 双対性により $\Ext^2(E,F)=\Hom(F,H^{-1}E)^*$ である。
右辺に非零射があれば、補題\ref{lem:rank-one-monomorphism}により単射であり、
余核 $T$ は $\mathcal A_1$ に属す。
その次数は式\eqref{eq:line-H}から
\[
 d(T)=d(H^{-1}E)-d(F)=d(E)-3-d(F)<0
\]
となり (EP) に矛盾する。
$E,F$ が同じ数値的線束類なら両向きの第二 Ext は零であり、
$\chi(E,F)=\chi(F,E)=1$ なので両向きに非零射が存在する。
いずれも単射であり、合成は非零である。
$\End(E)=\mathbb C$ だから合成は可逆なスカラーとなる。
一方の単射は左逆を持ち、他方も単射であることから互いに同型である。
最後に任意の線束対象の数値類は前系により $v_D$ なので、単射性が従う。
\end{proof}

\begin{remark}[存在・許容性・軌道を分ける]\label{rem:line-realization-set}
実現される因子類の集合を $\Omega_{\mathcal A}\subseteq\Lambda$ とすると、
上の仮定の下で $\operatorname{Line}(\mathcal A)\cong\Omega_{\mathcal A}$ である。
$H$ は線束対象を保つので $\Omega_{\mathcal A}-K=\Omega_{\mathcal A}$。
しかし $\Omega_{\mathcal A}=\Lambda$ は存在定理であり、一意性からは従わない。
また $\operatorname{Line}_{\rm adm}(\mathcal A)=\operatorname{Line}(\mathcal A)$ も
一般にはまだ示していない。
実際、線束対象 $E$ に対して $m\ge0$ なら
\[
 \dim\Hom(E,H^mE)
 =1+\tfrac32m(m+1)+\dim\Ext^1(E,H^mE),
\]
ここで第二 Ext の消滅は定理\ref{thm:intrinsic-line-criterion}から従う。
所定の切断数のためには交差 Ext の消滅、豊富性のためには生成条件が別に必要である。
さらに $H$ の軌道商と全 exact 自己同値の軌道商を同一視しない。
\end{remark}

\begin{proposition}[許容性の生成・消滅判定]\label{prop:line-admissibility-test}
定理\ref{thm:ac-serre-positivity}の数値仮定の下で、線束対象 $L$ が許容されることは、
次の二条件と同値である。
\begin{enumerate}
\item 有限和による全射 $\bigoplus_j H^{-a_j}L\twoheadrightarrow\OO_0$ が存在する。
\item $\Ext^{i}(L,H^mL)=0$ が $i>0$, $m\ge0$ に対して成立する。
\end{enumerate}
\end{proposition}
\begin{proof}
$L$ が許容なら第一条件は AZ 生成である。
第二 Ext は定理\ref{thm:numerical-line-uniqueness}で消え、
規定の切断数と Riemann--Roch の比較から第一 Ext も消える。
逆に二条件を仮定する。
第一条件を反標準 twist し、$\OO_0$ の任意 tail による生成と合成すると、
$L$ の任意 tail の有限和も全対象を生成する。
任意の $M$ に対して $P=\bigoplus_jH^{-b_j}L\twoheadrightarrow M$ を取り、
その核を $N$ とする。
負の自己 Hom 消滅と Serre 双対性から
$\Ext^2(L,H^nM)=\Ext^2(L,H^nN)=0$ が $n\gg0$ で成立する。
第二条件は $\Ext^1(L,H^nP)=0$ を与えるから、長完全列により
$\Ext^1(L,H^nM)=0$ も従う。
任意の全射の核にこの消滅を適用すれば AZ の Hom 全射条件を得る。
切断数は第二条件と Riemann--Roch、負方向は
定理\ref{thm:intrinsic-line-criterion}の負の自己 Hom 消滅から得られる。
従って $L$ は許容構造対象である。
\end{proof}

\subsection{零次数単純対象と許容性の障害}
ここでも定理\ref{thm:ac-serre-positivity}の数値仮定を置く。
次の命題は、線束対象であることから AZ 豊富性へ進む際に
調べるべき障害を与える。障害となる単純対象の実在は仮定であり、
本稿はその具体例を構成したとは主張しない。

\begin{lemma}[零次数単純対象の数値類]\label{lem:zero-degree-simple-classification}
$d(T)=0$ なる単純対象 $T$ に対して $C=c_1(T)$ と置く。
このとき $r(T)=0$ であり、$C=0$ または
\[
 C^2=-2,\qquad HT\cong T,\qquad \Ext^1(T,T)=0
\]
のいずれかである。後者では
$c=\dim\Hom(\OO_0,T)=\chi(\OO_0,T)>0$ となる。
\end{lemma}
\begin{proof}
$T\in\mathcal A_0\subset\mathcal A_1$ だから $r(T)=0$ である。
Hom 有限性と単純性により $\End(T)=\mathbb C$。
$H^{-1}T$ も単純なので
\[
 \epsilon:=\dim\Ext^2(T,T)=\dim\Hom(T,H^{-1}T)\in\{0,1\},
 \qquad \dim\Ext^1(T,T)=1+\epsilon+C^2.
\]
$C\in K^\perp$、$K^2=3$ と $\Lambda$ の符号 $(1,6)$ から
$K^\perp$ は負定値である。
また $D_0=c_1(\OO_0)$ とすると
$\chi(\OO_0,T)=s(T)-D_0\cdot C$ は整数だから $s(T)$ は整数であり、
$s(T)-C^2/2\in\mathbb Z$ から $C^2$ は偶数である。
従って $C\ne0$ なら $C^2=-2$、$\epsilon=1$、$\Ext^1(T,T)=0$ となる。
$\epsilon=1$ は $HT\cong T$ を意味する。
この場合、定理\ref{thm:ac-serre-positivity}と $HT\cong T$ から
$\Ext^{>0}(\OO_0,T)=0$ を得る。
同定理の零階数・零次数の場合の正値性により
$c=\chi(\OO_0,T)>0$ である。
\end{proof}

\begin{lemma}[数値的点への Ext 消滅]\label{lem:numerical-point-ext}
$P$ を $\operatorname{ch}(P)=(0,0,s)$ なる単純対象とする。
このとき $s>0$ であり、任意の線束対象 $L$ について
\[
 \Ext^1(P,L)=\Ext^1(L,P)=0,\qquad \dim\Hom(L,P)=s.
\]
\end{lemma}
\begin{proof}
$s=\chi(\OO_0,P)>0$ は定理\ref{thm:ac-serre-positivity}から従う。
非分裂な拡大 $0\to L\to F\to P\to0$ が存在すると仮定する。
$F$ の非零 torsion 部分対象は $L$ と交わらず $P$ に単射するため、
単純性から $P$ と同型になり、拡大を分裂させてしまう。
従って $F$ は rank-one torsion-free である。
しかし $\operatorname{ch}(F)=\operatorname{ch}(L)+(0,0,s)$ なので
$n(F)=-s<0$ となり、系\ref{cor:line-numerical-defect}に矛盾する。
これで $\Ext^1(P,L)=0$ を得る。
$H^{-1}L$ も線束対象であり、Serre 双対性から
$\Ext^1(L,P)\cong\Ext^1(P,H^{-1}L)^*=0$。
さらに $\Ext^2(L,P)=\Hom(P,H^{-1}L)^*=0$ なので、
$\chi(L,P)=s$ が Hom 次元を与える。
\end{proof}

\begin{theorem}[root 単純対象に対する線束の反射]\label{thm:root-line-reflection}
定理\ref{thm:ac-serre-positivity}の仮定を置く。
$T$ を $d(T)=0$, $C=c_1(T)\ne0$ なる単純対象とし、
$\operatorname{ch}(T)=(0,C,s)$ と書く。
任意の線束対象 $L$ に対して $D=c_1(L)$、
$a=\chi(L,T)=s-D\cdot C$ と置くと、
\begin{equation}\label{eq:root-line-hom-ext}
 \dim\Hom(L,T)=\max(a,0),\qquad
 \dim\Ext^1(L,T)=\dim\Ext^1(T,L)=\max(-a,0).
\end{equation}
線束の同型類集合上に involution $\rho_T$ があり、
\begin{equation}\label{eq:root-affine-reflection}
 c_1(\rho_TL)=D-aC=D+(D\cdot C-s)C
\end{equation}
を満たす。$a>0$ では評価全射の核、$a<0$ では普遍拡大、
$a=0$ では $L$ 自身によって実現される：
\[
 \begin{array}{ll}
 a>0:&0\longrightarrow\rho_TL\longrightarrow L
          \longrightarrow T^{\oplus a}\longrightarrow0,\\[2pt]
 a<0:&0\longrightarrow L\longrightarrow\rho_TL
          \longrightarrow T^{\oplus(-a)}\longrightarrow0.
 \end{array}
\]
さらに $R_{\rho_TL}\cong R_L$ である。
\end{theorem}
\begin{proof}
補題\ref{lem:zero-degree-simple-classification}により
$C^2=-2$, $HT\cong T$, $\Ext^1(T,T)=0$ である。
torsion-free 性から $\Hom(T,L)=0$ と
$\Ext^2(L,T)=\Hom(T,H^{-1}L)^*=0$ を得る。
Serre 双対性と $HT\cong T$ より
$\Ext^1(T,L)\cong\Ext^1(L,T)^*$。
$h=\dim\Hom(L,T)$、$b=\dim\Ext^1(T,L)$ と置けば $a=h-b$ である。

$V=\Ext^1(T,L)$ とし、
$\Ext^1(V\otimes T,L)=V^*\otimes V$ の恒等元による拡大
\[
 0\longrightarrow L\longrightarrow F\longrightarrow V\otimes T
 \longrightarrow0
\]
を取る。境界写像 $\Hom(T,V\otimes T)=V\to\Ext^1(T,L)=V$ は同型なので、
$\Hom(T,F)=0$ である。
$F$ の torsion 部分対象は $L$ と交わらず、半単純対象 $V\otimes T$ に埋め込まれる。
非零なら $T$ を部分対象として含み、上の消滅に反する。
従って $F$ は torsion-free であり、$qF\cong qL$ は非零単純なので階数一である。
数値的欠陥を計算すると
\[
 \begin{aligned}
 n(F)
 &=\frac{(D+bC)^2}{2}-\left(\frac{D^2}{2}+bs\right)\\
 &=-b(s-D\cdot C)-b^2=-b(h-b)-b^2=-bh.
 \end{aligned}
\]
系\ref{cor:line-numerical-defect}より $n(F)\ge0$ だから $bh=0$。
$h,b\ge0$ と $a=h-b$ から式\eqref{eq:root-line-hom-ext}が従う。

$a<0$ なら $b=-a$ であり、上の $F$ は $n(F)=0$ なので線束対象である。
$a>0$ なら $h=a$ であり、全 Hom による評価 $L\to T^{\oplus a}$ は全射である。
実際、像が真の部分対象なら、それを殺す非零射 $T^{\oplus a}\to T$ が
$\Hom(L,T)$ の基底間に非自明な関係を与える。
核 $E$ は rank-one torsion-free であり、
\[
 n(E)=a(s-D\cdot C)+\frac{a^2C^2}{2}=a^2-a^2=0
\]
なので線束対象である。これで各場合の $\rho_TL$ を得る。
どちらの場合も式\eqref{eq:root-affine-reflection}が成立し、
$s-(D-aC)\cdot C=-a$ である。
二回反射した Chern 類は $D$ に戻るので、
定理\ref{thm:numerical-line-uniqueness}から $\rho_T^2L\cong L$ を得る。
従って構成時の基底の選択も対象の同型類には影響しない。

環同型は $a>0$ の評価列で証明すればよい。
式\eqref{eq:root-line-hom-ext}より $\Ext^1(L,T)=0$ であり、
評価列から $\Hom(E,T)=0$ となる。
任意の整数 $n$ に対して
\[
 \Hom(T,H^nL)=0,\qquad
 \Ext^1(T,H^nL)\cong\Ext^1(L,H^{-n-1}T)^*=0.
\]
評価列に $\Hom(-,H^nL)$ を適用し、次に
$\Hom(E,H^nT)=0$ を使うと、
\[
 \Hom(L,H^nL)\cong\Hom(E,H^nL)\cong\Hom(E,H^nE)
\]
を得る。これは $L\to H^nL$ を $E$ に制限する写像であり、
$H$ と射の合成に両立するので、$R_L\cong R_E$ を与える。
$a<0$ の場合は反射後の正の側からこの議論を適用し、
$a=0$ の場合は恒等写像を取ればよい。
\end{proof}

\begin{corollary}[許容線束の半空間条件と root 排除]\label{cor:admissible-root-halfspace}
上の $T$ に対し、$L\in\operatorname{Line}_{\rm adm}(\mathcal A)$ なら
\[
 s-c_1(L)\cdot C>0.
\]
特にすべての $D\in\Lambda$ が許容線束の第一 Chern 類として実現されるなら、
零次数単純対象はすべて $c_1=0$ を満たす。
\end{corollary}
\begin{proof}
$a\le0$ なら定理より $\Hom(L,T)=0$ であり、$HT\cong T$ なので
$\Hom(H^{-n}L,T)=0$ がすべての $n$ で成立する。
これは $(L,H)$ の AZ 生成条件に反する。
最後の主張では、$C\ne0$ なる root simple があれば $D=-mC$ を取り、
$s-D\cdot C=s-2m\le0$ となる十分大きい整数 $m$ を選ぶと矛盾する。
\end{proof}

\begin{remark}
式\eqref{eq:root-affine-reflection}は $D\cdot C=s$ を壁とする affine reflection であり、
線束 $L$ 自体を変える。固定弱標識上の有限 Weyl 標識変更とは区別する。
構成は spherical twist の評価・普遍拡大と同じ数値的形式を持つが、
ここでは固定 heart 内の完全列によって証明した。
環同型から $L$ と $\rho_TL$ を交換する exact 自己同値が得られるとは主張しない。
半空間条件は必要条件であり、その正の側にある全線束の豊富性を結論しない。
\end{remark}

\begin{proposition}[同じ切断環を持つ非許容線束]\label{prop:root-simple-obstruction}
零次数 root 型単純対象 $T$ が存在すると仮定し、
$c=\dim\Hom(\OO_0,T)>0$ とする。
評価全射の核
\[
 0\longrightarrow E\longrightarrow\OO_0\longrightarrow T^{\oplus c}
 \longrightarrow0
\]
は線束対象であり、制限写像は次数付き環同型 $R_{\OO_0}\cong R_E$ を与える。
一方 $(E,H)$ は AZ 豊富ではない。
\end{proposition}
\begin{proof}
補題\ref{lem:zero-degree-simple-classification}と
定理\ref{thm:root-line-reflection}を $L=\OO_0$ に適用する。
反射後は $\chi(E,T)=-c$ なので $\Hom(E,T)=0$ であり、
$HT\cong T$ から全ての $n$ に対して $\Hom(H^{-n}E,T)=0$ となる。
従って AZ 生成条件が成立しない。
環同型は同定理の制限写像による証明で与えられる。
この命題は $T$ の存在を仮定するものであり、実例の構成ではない。
\end{proof}


\subsection{AS heart と反標準 heart の比較}
本稿で比較するのは $\mathcal A_{\rm AS}$ と $\mathcal A_{\rm ac}$ だけである。
次の十分条件は、指定した導来同値を exact 同値へ強める際に用いる。
\begin{lemma}[消滅と生成による heart の識別]\label{lem:ample-heart-transfer}
$\mathcal A,\mathcal B$ の有界導来圏の同値 $\Phi$ と、
各 heart 内の列 $P_n,Q_n$ ($n\ge0$) があるとする。
各対象 $M$ に対し十分大きな $n$ で
$\Ext^{p}(P_n,M)=0$ ($p>0$) かつ有限直和 $P_n^{\oplus a}\twoheadrightarrow M$ が存在し、
$Q_n$ についても同じ性質を仮定する。
十分大きなすべての $n$ で $\Phi(P_n)\cong Q_n$ なら $\Phi$ は $t$-exact である。
\end{lemma}
\begin{proof}
$C\in D^b(\mathcal A)$ の有限個の非零 cohomology 対象に対して
一様に大きい $n$ を取る。標準 truncation と Ext 消滅から
\[
 \Hom(P_n,C[i])\cong\Hom(P_n,H^i(C))
\]
となる。生成全射の存在によって右辺がすべての十分大きい $n$ で零であることと
$H^i(C)=0$ は同値である。
従って $C$ が $\mathcal A$ の heart に属すことは、
$i\ne0$ に対する上の Hom 消滅で特徴付けられる。
$\Phi$ とその逆に同じ議論を適用すれば、二つの heart が一致する。
\end{proof}


この補題では、両側の列をそれぞれ指定された構造対象の
十分負の反標準列として取る。両側での消滅・生成と列の同定は、
数値的な $K_0$ の一致や導来同値の存在だけでは代用できない。
AS heart で得た線束を反標準 heart の線束と認識する際も、
このような比較の仮定を明記する。

可換点を特殊ファイバーとする許容相対 Artin 変形では、
各古典的線束の一意な持上げは定理\ref{thm:formal-structure}で与えられる。
全数値類の同時存在は、後の定理\ref{thm:very-general-line-existence}で
明示する代数族の仮定の下では非常に一般の点で証明できる。
その適用には指定された反標準 heart と代数族との比較が必要であり、
任意の非可換点での存在や全域的な開性は主張しない。
